\section{Joint information conditional on base paths}\label{sec:base-paths}

A marginal estimate with no conditioning does not describe what remains
random after a prescribed set of paths has been observed. The next
theorem supplies that stronger information. Its reference law depends
on the observed endpoints, and its error is averaged over the actual
law of the observed paths. This is the form needed when a subsequent
construction reveals a color-occupancy history.

\begin{theorem}[A joint estimate given base paths]\label{thm:base-path-joint}
Fix $0<p<1$, and put
\[
 q=\frac{1+p}{2},\qquad \rho=\frac{1+q}{2}=\frac{3+p}{4},
 \qquad c_p=-\frac18\log_2\rho,\qquad
 b=b(p)=\left\lceil\frac{10}{c_p}\right\rceil.
\]
Run independent fair matching switches on $V=\mathbb F_2^d$ in any
fixed cyclic order of its $d$ coordinate directions, from any deterministic
starting deck. Let $n=2^d$, $T=2bd$, and fix a set $B$ of $r_0$ labels
and an ordered list $(a_1,\ldots,a_k)$ of other distinct labels, where
$r_0+k-1\le pn$. Write $\mathcal H_B$ for the full labeled paths of
$B$ through time $T$, and $A_B=V\setminus\Pi_T(B)$. For all sufficiently
large $d$ (depending only on $p$),
\begin{equation}\label{eq:base-path-joint}
 \E\left\|
 \mathcal L\big((\Pi_T(a_j))_{j=1}^k\mid\mathcal H_B\big)
 -\Unif\big(\operatorname{inj}([k],A_B)\big)
 \right\|_{\TV}\le k n^{-4}.
\end{equation}
The same assertion holds with $\mathcal H_B$ replaced by any coarser
sigma field that determines $A_B$. In particular it holds in the
reverse cyclic coordinate order, and hence for the inverse permutation
of a block of $2b$ complete sweeps, with the base paths observed in that
reversed process. The empty-list error is zero.
\end{theorem}

The proof first contracts a scalar vector uniformly over every
deterministic observed set. We then reveal the additional list one
card at a time, retaining the base paths throughout.
The next lemma also allows an arbitrary deterministic sum-zero input
vector, not just the centered point mass used in the first route.
The half-density estimate already controls the noise of such a vector.
Combining it with the one-step memory inequality contracts the energy
during every layer of the second sweep. The resulting bound can then
be restarted from each realized available set and vector.

\begin{lemma}[Uniform two-sweep contraction]\label{lem:base-path-energy}
Declare any $r\le pn$ labels observed. Let $f_0$ be a deterministic
real vector, supported on the available positions, with sum zero.
Propagate $f_0$ by the conditional transport-and-average rule given
the observed paths, and extend it by zero on occupied positions.
For all sufficiently large $d$,
\[
 \E\|f_{2d}\|_2^2\le n^{-c_p}\|f_0\|_2^2.
\]
Consequently, for every positive integer $a$,
\[
 \E\|f_{2ad}\|_2^2\le n^{-a c_p}\|f_0\|_2^2.
\]
All norms use counting measure on $V$.
\end{lemma}

\begin{proof}
Use the layer numbering of Section~\ref{sec:first-route}: layer $s$
takes $f_s$ to $f_{s+1}$. Put $a_s=\E\|f_s\|_2^2$ and let $w_s$
be the blocked energy in Lemma~\ref{noise:memory}. With
$\alpha=1-p$ and
$\eta_n=2\exp(-(1-p)^2n/4)$, Lemma~\ref{noise:halves} gives
\[
 w_s\le q a_s+\eta_n\|f_0\|_2^2\qquad(s\ge d).
\]
The identity \eqref{noise:one-step-memory} therefore yields
\begin{equation}\label{eq:base-noise}
 a_{s+1}\le\rho a_s+\tfrac12\eta_n\|f_0\|_2^2
 \qquad(s\ge d).
\end{equation}
Since $a_d\le\|f_0\|_2^2$, iterating over layers
$d,d+1,\ldots,2d-1$ gives
\[
 a_{2d}\le
 \left(\rho^d+\frac{\eta_n}{2(1-\rho)}\right)\|f_0\|_2^2
 \le n^{-c_p}\|f_0\|_2^2
\]
for sufficiently large $d$. Here $\rho^d=n^{-8c_p}$, and the
exponential error is smaller than every fixed inverse power of $n$.

At each later two-sweep boundary, condition on the observed paths only
up to that boundary. The available set and propagated vector are then
fixed, while future switch coins remain independent and fair. The
estimate is uniform in this set and vector, so its conditional version
iterates to give the second assertion. This conditioning on the past
is distinct from the full-path conditional law used to define a tagged
vector: prefix measurability identifies that vector, but the energy
estimate averages over the future observed paths.
\end{proof}

\begin{proof}[Proof of Theorem~\ref{thm:base-path-joint}]
For one unobserved tag, start with its point mass minus uniform measure
on the $n-r$ available positions. Its squared norm is $1-1/(n-r)\le1$.
The preceding lemma and $bc_p\ge10$ show that its endpoint conditional
law, given all $r$ observed paths, has expected variation distance at
most
\[
 \frac12\E\|f_T\|_1
 \le\frac12\sqrt{n\E\|f_T\|_2^2}
 \le\frac12 n^{-9/2}\le n^{-4}
\]
from uniform on the endpoint available set.

For the list, at step $j$ regard $B\cup\{a_1,\ldots,a_{j-1}\}$
as observed. Its size is at most $pn$, so the preceding expected
one-card bound applies to this enlarged observed set. Averaging over the paths of
$a_1,\ldots,a_{j-1}$ while keeping their endpoints and the base
paths can only decrease the expected variation distance: the comparison
law depends on those endpoints alone. For each fixed base history,
the sequential comparison lemma then bounds the joint error by the
sum of the one-card errors, averaged under the actual prefix law.
One may see this directly by replacing the last sampling kernel by
uniform sampling from the available set, and then replacing earlier
kernels in reverse order. Each replacement costs its expected
conditional error; all kernels appended after it are contractions
in total variation. Averaging also over the base history gives
\eqref{eq:base-path-joint}.

If a coarser sigma field determines the endpoint available set, the
same uniform reference is constant on each of its atoms. Conditional
convexity therefore gives the asserted coarsening. No step in the
scalar argument depended on which cyclic order was chosen. Finally,
the inverse of a switch block reverses its independent involutive
layers. Starting that reversed sequence from a deterministic deck gives
the inverse increment law, with its own observed base paths, and proves
the reverse-order conclusion.
\end{proof}
